\section{Modular transformations and prime-filtered cusp errors}
\label{sec:modular-case}
\record{M01}\record{M02}
A modular transformation can turn slow convergence near one boundary
point into fast convergence near another. For this problem it also
introduces an unwanted multiple of $\pi^4$. The construction below
removes that term exactly. It then separates three questions that are
easy to confuse: whether the transformed error is small and nonzero,
whether the coefficients have manageable denominators, and whether
these facts produce small integer forms.

\subsection{The classical identity and its cancellation}
Write $\chi=\chi_4$, $q=e^{2\pi i\tau}$, with $\Im\tau>0$, and define
\begin{align}
 L_F(\tau)&=\sum_{d\geq1}\frac{\chi(d)q^d}{d^4(1-q^d)},&
 L_G(\tau)&=\sum_{d\geq1}\frac{q^d}{d^4(1+q^{2d})},\notag\\
 F(\tau)&=L_F(\tau)-2L_F(2\tau)
       =\sum_{d\text{ odd}}\frac{\chi(d)q^d}{d^4(1+q^d)},&
 G(\tau)&=L_G(\tau)-\frac{L_G(2\tau)}{16}
       =\sum_{d\text{ odd}}\frac{q^d}{d^4(1+q^{2d})}.
 \label{eq:modular-series}
\end{align}
All sums converge absolutely. A specialization of Bradley's
Ramanujan-type formula~\cite[Corollary 4]{Bradley} gives
\begin{equation}\label{eq:modular-LF}
 L_F(\tau)+\frac\theta2
 =\frac{i\tau^3}{2}L_G\left(-\frac1{4\tau}\right)
   +i\pi^4\left(\frac5{3072\tau}-\frac\tau{192}
                                  -\frac{\tau^3}{360}\right).
\end{equation}
For clarity about normalization, in that corollary take the integer
parameter $2$, put its positive parameters equal to
$\alpha=\pi\Im\tau$, $\beta=\pi/\Im\tau$, and use
$E_0=1,E_2=-1,E_4=5$, $B_0=1,B_2=1/6,B_4=-1/30$.
This gives~\eqref{eq:modular-LF} on the positive imaginary axis;
analytic continuation gives the upper half-plane identity.
Subtract twice the same formula at $2u$ from the formula at $u$,
then put $u=-1/(8\tau)$. The result is
\begin{equation}\label{eq:modular-fricke}
 \tau^3F\left(-\frac1{8\tau}\right)-\frac{i}{64}G(\tau)
 =\frac\theta2\tau^3-i\pi^4
                \left(\frac{\tau^2}{512}+\frac1{12288}\right).
\end{equation}
Thus the special value enters a transformation law with an explicit
polynomial correction. This is the Eichler-integral viewpoint:
four derivatives $q\,d/dq$ of these series are holomorphic modular
forms of weight five. Their Fourier expansions and transformations
are instances of the Eisenstein-series construction in
Berndt--Straub~\cite[Section 4 and Example 6.4]{BerndtStraub}.

Choose coefficients
\[
 a_1=16,\qquad a_2=-10,\qquad a_4=1.
\]
They satisfy $\sum a_rr=\sum a_rr^3=0$ and $\sum a_r=7$.
Consequently, with
\begin{equation}\label{eq:modular-AH}
 A=16F-10F(2\tau)+F(4\tau),\qquad
 H=A+\frac{\sqrt2}{2}\bigl(4G-5G(2\tau)+G(4\tau)\bigr),
\end{equation}
all $\pi^4$ terms cancel and
\begin{equation}\label{eq:modular-H}
 H\left(-\frac1{32\tau}\right)-\frac72\theta
  =\frac{i}{32^{3/2}\tau^3}\left(H(\tau)-\frac72\theta\right).
\end{equation}
Here is a direct way to check the coefficients, without guessing a
modular combination. Solve~\eqref{eq:modular-fricke} for $F(r\tau)$.
Its polynomial terms are proportional to $r\tau$ and $r^3\tau^3$,
and its $G$ term is
$-8i(r\tau)^3G(-1/(8r\tau))$. Under $\tau\mapsto-1/(32\tau)$
the indices $r=1,2,4$ reverse to $4/r$; their cubic weights give
exactly $(\sqrt2/2)(4,-5,1)$. The two displayed zero sums remove
the polynomial terms in both directions. At the fixed point
$\tau=i/\sqrt{32}$ the multiplier in~\eqref{eq:modular-H} is $-1$,
so $H(i/\sqrt{32})=7\theta/2$.

\subsection{Deleting prime denominators while retaining an exact error}
Let $S$ be a finite set of odd primes, $P=\prod_{p\in S}p$, and
$\alpha_S=\prod_{p\in S}(1-\chi(p)/p^4)$. Empty products are one.
For $V_df(\tau)=f(d\tau)$ set
\begin{equation}\label{eq:modular-filter}
 A_S=\prod_{p\in S}(1-\chi(p)p^{-4}V_p)A
    =\sum_{r\in\{1,2,4\}}\sum_{d\mid P}
           a_r\mu(d)\chi(d)d^{-4}F(rd\tau),\qquad
 c_S=\frac72\alpha_S.
\end{equation}
The symbol $\mu(d)$ here is $(-1)$ to the number of prime factors of
the squarefree divisor $d$. The elementary identity
\[
 (1-\chi(p)p^{-4}V_p)F
   =\sum_{\substack{d\text{ odd}\\p\nmid d}}
                  \frac{\chi(d)q^d}{d^4(1+q^d)}
\]
shows that every coefficient of $A_S$ is $p$-integral for $p\in S$.
It removes an outer divisor in the Lambert sum; it does not remove
every Fourier coefficient whose index is divisible by $p$.
Dominated convergence gives $A_S(iy)\to c_S\theta$ as $y\downarrow0$.

Apply the solved form of~\eqref{eq:modular-fricke} term by term in
\eqref{eq:modular-filter}. The same two zero sums eliminate the
polynomial terms for every $d$, giving the exact remainder
\begin{equation}\label{eq:modular-cusp-exact}
 A_S(\tau)-c_S\theta
   =-8i\tau^3\sum_{r,d\mid P}
          \frac{a_rr^3\mu(d)\chi(d)}{d}
                         G\left(-\frac1{8rd\tau}\right),
\end{equation}
where the sum over $r$ is still $1,2,4$.
This identity yields a uniform sign and size estimate.

\begin{proposition}\label{prop:modular-cusp}
Put $\sigma_P=\sum_{d\mid P}d^{-1}$, $x=\pi/(16Py)$ and $z=e^{-x}$.
If $x\geq\log(20P\sigma_P)$, then
\begin{equation}\label{eq:modular-cusp-bound}
 A_S(iy)-c_S\theta
  =-\frac{512\mu(P)\chi(P)y^3z}{P}(1+\varepsilon),\qquad
 |\varepsilon|\leq2z^2+5P\sigma_Pz<\frac12.
\end{equation}
In particular the error is nonzero, with the displayed sign, and its
absolute value lies strictly between $256y^3z/P$ and $768y^3z/P$.
\end{proposition}
\begin{proof}
Regard $G$ temporarily as a power series in a real argument $z$.
For $0<z\leq1/2$ its defining sum gives
$|G(z)-z|\leq2z^3$ and $0<G(z)\leq2z$:
the first summand differs from $z$ by at most $z^3$, while the
remaining odd summands are bounded by
$z^3/[3^4(1-z^2)]$.
In~\eqref{eq:modular-cusp-exact} the unique largest dilation is
$rd=4P$, with coefficient $64\mu(P)\chi(P)/P$ and argument $z$.
Every other argument is $z^{4P/(rd)}$ with integer exponent at least
$2$. Since $\sum_r|a_r|r^3=160$, their total absolute contribution
inside the sum is at most $320\sigma_Pz^2$. Divide by $64z/P$ and
include the leading term's relative error $2z^2$.
Finally $z\leq(20P\sigma_P)^{-1}\leq1/20$ makes the bound less
than $1/2$.
\end{proof}
The required approach scale is explicit:
$y\leq\pi/[16P\log(20P\sigma_P)]$. These are analytic evaluations
at $q=e^{-2\pi y}$. They are not rational approximants merely because
their error is small. A further construction must turn the identity
into rational coefficients without losing its decay.

\subsection{A fixed-filter obstruction to the ordinary contour estimate}
We give the analytic limitation rather than infer it from a plot of
singularities. A holomorphic modular form of positive integer weight,
periodic with period one and holomorphic at the cusps of a finite-index
subgroup, cannot extend meromorphically through a finite rational cusp
unless it is zero. Indeed, choose
$\gamma=\left(\begin{smallmatrix}a&b\\c&d\end{smallmatrix}\right)$
with $c\ne0$ taking infinity to the cusp. After increasing the cusp
width if necessary, $f|_k\gamma=a_0+O(e^{-\eta\Im z})$ there.
Meromorphic continuation through $a/c$ would make this a convergent
Laurent series in $1/z$. Its difference from $a_0$ is exponentially
flat on the imaginary axis, so every Laurent coefficient is zero.
It follows that $f(\tau)=a_0(a-c\tau)^{-k}$, which contradicts
periodicity if $a_0\ne0$ and nonzeroness if $a_0=0$.

For fixed $S$, $(q\,d/dq)^4A_S$ is such a modular form after passing
to a common subgroup for its finitely many dilations. Its coefficient
of $q$ is $16$, so it is nonzero. If $A_S$ continued through a root
of unity, its fourth derivative would too. Every root of unity is
therefore a singularity; their density makes the unit circle a
\emph{natural boundary}, across which no analytic continuation is possible.

There are two useful consequences. First, for $h=2A_S$, the cumulative
germ $(X-h(q))/(1-q)$ is nonrational for every complex $X$. The
Hankel criterion used in Section~\ref{sec:common-base} therefore gives
infinitely many nonzero endpoint determinants for each fixed filter.
Second, the compact set removed in the ordinary reciprocal-plane
absolute-contour estimate must contain the whole unit circle, whose
logarithmic capacity is one. The estimate used in
Section~\ref{sec:common-base} cannot supply a negative quadratic decay
constant here. This is a limitation of that estimate. It is neither
a lower bound for the actual determinants nor a uniform statement
when $S$ grows with the matrix size.

\subsection{An explicit coefficient realization and its prime cost}
For readers wishing to reproduce the coefficient experiment, the
level-32 choice was the following formal change of variable and prefactor:
\begin{align*}
 t(q)&=q\prod_{d\in\{1,2,4,8,16,32\}}\prod_{j\geq1}
                 (1-q^{dj})^{e_d},
       &(e_1,e_2,e_4,e_8,e_{16},e_{32})&=(2,-3,1,1,-3,2),\\
 J(q)&=\vartheta_3(2\tau)^2+2\sqrt2\,\vartheta_3(16\tau)^2,
       &\vartheta_3(u)&=\sum_{m\in\Z}e^{\pi i m^2u},\\
 Y(q)&=q\,\frac{d}{dq}\log t(q),&
 \mathcal A(t(q))&=\frac{J(q)^5}{Y(q)},\qquad
 \mathcal B(t(q))=\mathcal A(t(q))H(\tau).
\end{align*}
These definitions alone suffice to compute every coefficient exactly:
$t=q+O(q^2)$ has an integral compositional inverse,
$Y\in1+q\Z[[q]]$, and all operations take place in
$\Q(\sqrt2)[[q]]$. Write
$\mathcal A=\sum(a_{n,0}+\sqrt2a_{n,1})t^n$ and
$\mathcal B=\sum(b_{n,0}+\sqrt2b_{n,1})t^n$.
Then $a_{n,j}\in\Z$, and $2\operatorname{lcm}(1,\ldots,n)^4b_{n,j}$
is integral. The latter follows directly from the Lambert coefficients
in~\eqref{eq:modular-series}; integral composition and multiplication
preserve that denominator bound.

The fourth powers are real costs of these coefficient sequences, not
just artifacts of this sufficient bound. For an odd prime $p$, all
coefficients before degree $p$ are $p$-integral, while
\[
 p^4[q^p]H\equiv16\chi(p)+2\sqrt2\pmod p.
\]
Since $\mathcal A(0)=(1+2\sqrt2)^5=401+298\sqrt2$, and composition
with $q=t+O(t^2)$ does not change this first nonintegral coefficient,
\begin{equation}\label{eq:modular-prime-coeff}
 \begin{split}
 p^4b_{p,0}&\equiv6416\chi(p)+1192\pmod p,\\
 p^4b_{p,1}&\equiv4768\chi(p)+802\pmod p.
 \end{split}
\end{equation}
Thus $v_p(b_{p,j})=-4$ apart from the finitely many primes dividing
one of the four displayed nonzero integers. The least denominator
clearing a whole prefix through $n$ consequently has logarithm
$4n+o(n)$, by the prime number theorem and the matching upper bound.
This does not say that every individual coefficient has that denominator,
or determine the primitive content of a determinant built from them.

\paragraph{What remains.}
The transformation and cusp estimate are available exactly, and the
coefficient construction is reproducible. A continuation of this route
must control primitive coefficient combinations, or supply a different
arithmetic interpretation of the transformed values. Neither selected
prime deletion nor the nonzero cusp error supplies this missing step.
The more specialized coordinate and rational-system tests remain in
records \recordid{M03}--\recordid{M04}; their finite comparisons do not
exclude other modular choices.
